\par\noindent\textbf{(1)} 设复透射振幅 t(ε)=t\_bg(ε+q)/(ε+i)，其中q为实Fano参数，这可拆成 t\_bg[1+(q−i)/(ε+i)]。第一项为非共振直接通道，第二项为带极点的离散共振响应；相干复振幅相加产生不同于Lorentz线的线型。
\par\noindent\textbf{(2)} 相对背景透射强度为 F(ε)=|t/t\_bg|²=(ε+q)²/(ε²+1)。分子在 ε=−q 恰为零，出现完全反共振（在无额外损耗且模型允许的条件下）；不能把两个通道的概率相加，否则得不到该零点。
\par\noindent\textbf{(3)} 求导 F'=2(ε+q)(1−qε)/(ε²+1)²。除反共振零点外，若q≠0另有极值 ε\_peak=1/q，其峰值F=1+q²。频率差 ω\_peak−ω\_zero=(Γ/2)(q+1/q)，解释线型的不对称程度并不只由Γ决定。
\par\noindent\textbf{(4)} 当|q|≫1并在 ε=O(1) 的窗口重新归一化，线型趋近Lorentz峰；当q=0，F=ε²/(1+ε²)为对称透明窗口。吸收可使q复数、反共振极小值非零；共振相位在近零点迅速变化，但带损情形不一定精确跳π。
